% !TeX root = ../cosmology_summary.tex


%% --------------------------------------------------------
\chapter{Ефект Закса-Вольфа та формування спектра CMB}
\label{sec:sachs}
%% --------------------------------------------------------

Кутова анізотропія температури реліктового випромінювання (CMB), що спостерігається нами сьогодні, є безпосереднім відбитком («знімком») стану акустичних осциляцій фотон-баріонної плазми в епоху рекомбінації Всесвіту ($z_{\mathrm{rec}} \approx 1100$). Флуктуації температури у заданому напрямку на небі $\hat{\mathbf{n}}$ розкладаються за сферичними гармоніками, а їхні статистичні властивості повністю описуються кутовим спектром потужності $\mathcal{D}_\ell^{TT} \equiv \frac{\ell(\ell+1)}{2\pi}C_\ell^{TT}$, де мультиполь $\ell$ обернено пропорційний кутовому масштабу $\theta \sim 180^\circ / \ell$ \cite{dodelson_2020}.

Спостережувані варіації температури $\delta T / T$ формуються внаслідок дії трьох основних фізичних механізмів на поверхні останнього розсіяння (LSS) \cite{hu_sugiyama_1995}:
\begin{equation}\label{eq:cmb_sources}
  \frac{\delta T}{\text{T}}(\hat{\mathbf{n}}) = \left( \frac{1}{4}\delta_\gamma + \Phi \right)_{\mathrm{LSS}} + \left( \hat{\mathbf{n}} \cdot \mathbf{v}_b \right)_{\mathrm{LSS}} + \int_{\eta_{\mathrm{LSS}}}^{\eta_0} \left( \Phi' + \Psi' \right) d\eta.
\end{equation}

Кожен із доданків у виразі (\ref{eq:cmb_sources}) відповідає за свій масштаб на спектрі:
\begin{enumerate}
  \item \textbf{Звичайний ефект Закса--Вольфа ($\frac{1}{4}\delta_\gamma + \Phi$):}
        Домінує на великих кутових масштабах (малі мультиполі $\ell \lesssim 100$),
        які на момент рекомбінації ще перебували поза горизонтом Хаббла.
        Він поєднує в собі власну термодинамічну анізотропію фотонів та
        гравітаційне червоне зміщення, якого фотони зазнають при виході
        з потенціальних ям $\Phi$.

        \textit{Адіабатичні збурення і умова $\frac{1}{4}\delta_\gamma = -\frac{2}{3}\Phi$.}
        Збурення називаються \emph{адіабатичними}, якщо всі компоненти
        (фотони, баріони, темна матерія) збурені узгоджено: локальний
        стан плазми відрізняється від середнього лише температурою, а
        не складом. Формально це означає, що відношення числових густин
        зберігається: $\delta(n_b/n_\gamma) = 0$. Оскільки
        $n_\gamma \propto T^3$ і $n_b \propto a^{-3}$, умова адіабатичності
        зводиться до $\delta_\gamma = \frac{4}{3}\delta_b$.

        На суперхабблових масштабах у радіаційно-домінованій епосі
        рівняння збереження для фотонів і метричний потенціал $\Phi$
        пов'язані через рівняння Пуассона:
        $k^2\Phi \approx -4\pi G a^2 \bar\rho\,\delta$.
        У граничному випадку $k \to 0$ (масштаби значно більші за горизонт)
        розв'язок для адіабатичних початкових умов дає
        \cite{hu_sugiyama_1995}:
        \begin{equation*}
          \frac{1}{4}\delta_\gamma = -\frac{2}{3}\Phi,
        \end{equation*}
        де мінус означає, що надлишок густини фотонів ($\delta_\gamma > 0$)
        відповідає потенціальній ямі ($\Phi < 0$): плазма стиснута саме
        там, де є гравітаційний мінімум. Підставляючи у вираз для ефекту
        Закса--Вольфа:
        \begin{equation*}
          \frac{1}{4}\delta_\gamma + \Phi
          = -\frac{2}{3}\Phi + \Phi = \frac{1}{3}\Phi.
        \end{equation*}
        Знак «$+$»: фотони виходять з ями ($\Phi < 0$) і зазнають
        гравітаційного червоного зміщення, тобто виглядають
        \emph{холоднішими} (фотон витрачає свою енергію (частоту) на подолання гравітаційного потенціалу ями, «червоніє», і для зовнішнього спостерігача ця область космосу здається плямою зниженої температури.). Ефект $\frac{1}{3}\Phi < 0$ означає
        дефіцит температури у напрямку потенціальної ями~---
        саме це формує плато Закса--Вольфа на великих масштабах.
  \item \textbf{Доплерівський ефект ($\hat{\mathbf{n}} \cdot \mathbf{v}_b$):} Зумовлений кінематичним рухом (швидкістю $\mathbf{v}_b$) фотон-баріонної плазми у полі збурень. Цей внесок зміщений за фазою відносно ефекту Закса--Вольфа і частково заповнює мінімуми між акустичними піками.
  \item \textbf{Інтегральний ефект Закса--Вольфа (ISW, інтегральний член):} Виникає, коли фотони CMB по дорозі до спостерігача пролітають крізь гравітаційні потенціали, що змінюються у часі ($\Phi' \neq 0$). Ранній ISW проявляється відразу після рекомбінації через залишковий вплив радіації, а пізній ISW~--- на надвеликих масштабах ($\ell < 10$) в сучасну епоху через прискорене розширення Всесвіту під дією темної енергії.
\end{enumerate}

%% --------------------------------------------------------
\section{Акустичні піки та затухання Сілка}
%% --------------------------------------------------------

На субхабблових масштабах ($\ell \gtrsim 100$) когерентні звукові коливання формують чітку серію \emph{акустичних піків} (рис.~\ref{fig:cmb_tt}). Перший і найвищий пік при $\ell \approx 220$ ($\theta \sim 1^\circ$) відповідає моді, яка встигла зазнати рівно одного максимального стиснення у потенціальній ямі до моменту рекомбінації. Його положення є чутливим індикатором просторової геометрії Всесвіту \cite{planck_2018}. На мультиполях $\ell \gtrsim 1000$ спектр експоненціально загасає через ефект \emph{дифузійного затухання} (затухання Сілка)~--- фотони мають скінченну довжину вільного пробігу і в процесі дифузії замивають дрібномасштабні флуктуації температури. Характерний масштаб затухання $\ell_D \sim 1000$--$1500$ визначається довжиною дифузії фотонів до рекомбінації \cite{silk_1968}.

\begin{figure}[t!]
\centering
\begin{tikzpicture}
\begin{semilogxaxis}[
    width  = 14.5cm,
    height = 8.5cm,
    xlabel = {Мультипольний момент $\ell$},
    ylabel = {$\mathcal{D}_\ell^{TT}\ [\mu\mathrm{K}^2]$},
    xmin   = 2,   xmax = 2600,
    ymin   = 0,   ymax = 6800,

    % --- РУЧНЕ КЕРУВАННЯ ДЛЯ УНИКНЕННЯ НАКЛАДАННЯ ---
    % Залишаємо лише ті тики, які реально підписати на такій шкалі без каші
    xtick = {2, 10, 30, 100, 500, 1000, 1500, 2000, 2500},

    % Форматуємо підписи вручну через макроси, щоб вони не налізали:
    xticklabels = {2, 10, 30, 100, 500, 1000, 1500, 2000, 2500},

    % АБО якщо потрібен СУВОРИЙ НАУКОВИЙ ФОРМАТ (sci) для великих чисел:
    % xticklabels = {2, 10, 30, 100, 500, $10^3$, $1.5{\cdot}10^3$, $2{\cdot}10^3$, $2.5{\cdot}10^3$},

    % Допоміжні тики (маленькі палички між цифрами), які роблять сітку «густою» як у книзі
    minor xtick = {3,4,5,6,7,8,9,20,40,50,60,70,80,90,200,300,400,600,700,800,900,1100,1200,1300,1400,1600,1700,1800,1900,2100,2200,2300,2400},

    ytick   = {0, 1000, 2000, 3000, 4000, 5000, 6000},
    grid   = both,
    grid style = {line width=0.3pt, draw=gray!25},
    major grid style = {line width=0.5pt, draw=gray!40},
    tick label style = {font=\small},
    label style      = {font=\small},
    legend style = {
        at={(0.6,0.97)}, anchor=north east,
        font=\small, fill=white, fill opacity=0.9,
        draw=gray!60, inner sep=4pt, row sep=1pt
    },
    clip = true,
]

%--- Theory curve (LCDM, CAMB) ---
\addplot[
    color     = theoryred,
    line width = 1.6pt,
    smooth,
] table[x=ell, y=Dl] {\theorydata};
\addlegendentry{$\Lambda$CDM (CAMB, Planck\,2018 best-fit)}

%--- Planck 2018 data with error bars ---
\addplot[
    only marks,
    mark       = *,
    mark size  = 1.4pt,
    color      = planckblue,
    error bars/.cd,
        y dir  = both,
        y explicit,
] table[
    x    = ell,
    y    = Dl,
    y error minus = errm,
    y error plus  = errp,
] {\planckdata};
\addlegendentry{Planck\,2018 (збіновані дані TT)}

%--- Annotate acoustic peaks ---
\node[font=\tiny, text=gray!70, above] at (axis cs:220,5732) {1-й пік};
\node[font=\tiny, text=gray!70, above] at (axis cs:540,2592) {2-й пік};
\node[font=\tiny, text=gray!70, above] at (axis cs:810,2540) {3-й пік};

%--- Sachs-Wolfe plateau annotation ---
\draw[->, gray!60, thin] (axis cs:6,1300) -- (axis cs:30,1070)
    node[midway, left, font=\tiny, text=gray!60] {плато Закса--Вольфа};

\end{semilogxaxis}
\end{tikzpicture}
\caption{Кутовий TT-спектр потужності CMB. Точки~--- дані Planck~2018
         \cite{planck_2018}.
         Крива~--- базова модель $\Lambda$CDM, розрахована CAMB
         з параметрами місії Planck~2018 \cite{planck_2018}
         ($H_0=67.36$, $\Omega_b h^2=0.02237$, $\Omega_c h^2=0.1200$,
         $\tau=0.0544$, $A_s=2.1\times10^{-9}$, $n_s=0.9649$).
         Видно три акустичні піки та плато Закса--Вольфа на великих
         кутових масштабах ($\ell \lesssim 30$).}
\label{fig:cmb_tt}
\end{figure}
